\providecommand{\devnum}[1]{#1}

% The preamble is expected to load amsthm.  The declarations below are made only
% if main.tex has not already made them.
\makeatletter
\@ifundefined{theorem}{%
  \theoremstyle{plain}%
  \newtheorem{theorem}{Theorem}[section]%
  \newtheorem{lemma}[theorem]{Lemma}%
  \newtheorem{corollary}[theorem]{Corollary}%
  \newtheorem{proposition}[theorem]{Proposition}%
  \theoremstyle{definition}%
  \newtheorem{definition}[theorem]{Definition}%
  \theoremstyle{remark}%
  \newtheorem{remark}[theorem]{Remark}%
}{}
\makeatother

\section{Theory}
\label{sec:theory}

Everything below is pathwise: one arrival sequence, one vector of service times,
one realisation of whatever randomness the base policy uses. No stationarity, no
independence and no assumption about score quality enters any proof, and the
upper bounds hold against a base policy that knows the future. An identity ties a
job's delay relative to first-come first-served to the work that overtook it; two
bounds turn that identity into the guarantee of Algorithm~\ref{alg:guard}; and a
set of constructions shows that each constant in those bounds is an exact
supremum. The proof of the identity is given here in full and the idea of every other
argument with it; Appendix~\ref{app:proofs} states the remaining ideas, and
Supplementary Sections~S4 and~S8 give every deferred proof and every extremal
construction.

\subsection{Conventions}

We use the model of Section~\ref{sec:model} throughout, and fix three readings
the proofs depend on. The dispatches at one instant form a \emph{phase}, ordered
as the policy performs them, and the dispatch sequence behind $j \prec i$ is the
concatenation of the phases; first-come first-served dispatches in increasing
rank inside a phase. A job dispatched in the same phase as $i$ executes no work
during $[a_i, s^P_i)$, that interval ending at the instant of the phase, yet
contributes its whole $C_j$ to $\mathrm{In}_i$. Remaining work at $a_i$ is read
\emph{after} the arrivals at $a_i$, so a job of rank below $i$ arriving at
exactly $a_i$ counts with its full service time.

\subsection{The workload gap between two policies}

Write $U_Q(t)$ for the unfinished work at time $t$ under policy $Q$, counting
every job with $a_j \le t$, and $N_Q(t)$ for the number of jobs present. The
unfinished-work vector of a $k$-server queue obeys the Kiefer--Wolfowitz
recursion~\cite{kiefer1955queues}, and first-come first-served minimises that
vector in the weak-majorization order for systems in which arrivals are
allocated to servers in a manner that does not depend on their service
times~\cite{daley1987fcfs}. The policies compared below are not restricted that
way, so what carries the
single-server argument to $k > 1$ here is a two-sided bound between two arbitrary
policies rather than an ordering.

\begin{lemma}[Workload comparison]
\label{lem:gap}
Let $Q_1$ and $Q_2$ be non-preemptive work-conserving $k$-server policies run on
the same input. Then $|U_{Q_1}(t) - U_{Q_2}(t)| \le (k-1)L$ for every $t$. For
$k \ge 2$ the inequality is strict, $|U_{Q_1}(t) - U_{Q_2}(t)| < (k-1)L$, and
the constant is the exact supremum, approached but never attained. At $k = 1$
both sides vanish: one server executes work at rate $1$ whenever any job is
present, so every work-conserving policy has the same unfinished work at every
instant.
\end{lemma}

The difference $D = U_{Q_1} - U_{Q_2}$ is continuous and piecewise linear, and
it cannot increase at any instant at which $U_{Q_1}$ exceeds $(k-1)L$, since
more than $k-1$ jobs are then present under $Q_1$ and every server of $Q_1$ is
busy; tracking $D$ back to the last instant at which $U_{Q_1}$ was below the
threshold gives the bound.
The proof, the strictness, and the cascade that shows the constant cannot be
lowered are in Appendices~\ref{app:strict} and~\ref{app:cascade}.

\subsection{The net-overtake identity}

For a job $i$ and a policy $Q$, let $R^Q_i$ be the remaining work at time $a_i$,
under $Q$, of the jobs of rank below $i$, read after the arrivals at $a_i$; and
let $\rho^Q_i$ be the remaining work at $s^Q_i$ of the jobs dispatched strictly
before $i$ in $Q$'s dispatch sequence that are still in service then. Write
\[
  \Delta_i \;=\; k\bigl(W_P[i] - W_{\mathrm{FCFS}}[i]\bigr)
              \;-\; \bigl(\mathrm{In}_i - \mathrm{Out}_i\bigr).
\]

\begin{theorem}[Net-overtake identity]
\label{thm:identity}
For every non-preemptive work-conserving $k$-server policy $P$, every input and
every job $i$,
\[
  \Delta_i \;=\; \bigl(R^P_i - R^{\mathrm{FCFS}}_i\bigr) - \rho^P_i
                 + \rho^{\mathrm{FCFS}}_i ,
  \qquad
  \bigl|\Delta_i\bigr| \;\le\; 2(k-1)L ,
\]
with strict inequality for $k \ge 2$; equivalently
$\bigl|\,W_P[i] - W_{\mathrm{FCFS}}[i] - (\mathrm{In}_i - \mathrm{Out}_i)/k\,\bigr|
 \le (2 - 2/k)L$. At $k = 1$ the error term vanishes and
$W_P[i] - W_{\mathrm{FCFS}}[i] = \mathrm{In}_i - \mathrm{Out}_i$ exactly.
\end{theorem}

\begin{proof}
\emph{Busy-period accounting.} Job $i$ waits throughout $[a_i, s^P_i)$, so by
work conservation every server is busy on that interval and exactly $k\,W_P[i]$
units of work are executed there, none of it on $i$. (If $s^P_i = a_i$ the
interval is empty and both sides vanish.) Split that work by rank.

A job $j$ of rank above $i$ has $a_j \ge a_i$, so it executes inside the window
only if it is dispatched inside it, that is only if $j \prec i$; conversely every
such $j$ is dispatched at or before $s^P_i$. The work executed on jobs of rank
above $i$ is therefore $\mathrm{In}_i - \rho^{\mathrm{new}}$, where
$\rho^{\mathrm{new}}$ is the remaining work at $s^P_i$ of those overtakers not
finished by then; a same-phase overtaker contributes $C_j$ to both terms and $0$
to the executed work, as it must.

Every job of rank below $i$ is present at $a_i$, with total remaining $R^P_i$.
At $s^P_i$ their total remaining is $\mathrm{Out}_i$, the jobs not dispatched
before $i$ each carrying their full service time, plus $\rho^{\mathrm{old}}$,
the unfinished part of those dispatched strictly before $i$, so the work
executed on them is $R^P_i - \mathrm{Out}_i - \rho^{\mathrm{old}}$. Adding the
two parts,
\begin{equation}
  k\,W_P[i] \;=\; R^P_i + \mathrm{In}_i - \mathrm{Out}_i - \rho^P_i ,
  \qquad \rho^P_i := \rho^{\mathrm{old}} + \rho^{\mathrm{new}} .
  \label{eq:busy}
\end{equation}

\emph{Bounding $\rho$.} The jobs counted in $\rho^P_i$ occupy distinct servers,
and none of them occupies the server that takes $i$, so there are at most $k-1$
of them and each has remaining work at most $L$: $0 \le \rho^P_i \le (k-1)L$.

\emph{The reference policy.} Under first-come first-served,
$\mathrm{In}_i = 0$, since a job of rank above $i$ dispatched before $i$ would
have to be chosen at an instant at which $i$ waits, and
$\mathrm{Out}_i = 0$, since at $s^{\mathrm{FCFS}}_i$ every job of rank below $i$
has already been dispatched. So \eqref{eq:busy} reads
$k\,W_{\mathrm{FCFS}}[i] = R^{\mathrm{FCFS}}_i - \rho^{\mathrm{FCFS}}_i$ with
$0 \le \rho^{\mathrm{FCFS}}_i \le (k-1)L$.

\emph{Comparison.} At $a_i$, after the arrivals there are processed, the two
systems contain the same jobs, and those of rank $i$ or above are untouched in
both, so $U_P(a_i) - U_{\mathrm{FCFS}}(a_i) = R^P_i - R^{\mathrm{FCFS}}_i$ and
Lemma~\ref{lem:gap} gives $|R^P_i - R^{\mathrm{FCFS}}_i| \le (k-1)L$, strictly
for $k \ge 2$. Subtracting
the two instances of \eqref{eq:busy} gives the identity. Each direction uses
only two of the three ranges: upward
$R^P_i - R^{\mathrm{FCFS}}_i \le (k-1)L$ and $\rho^{\mathrm{FCFS}}_i \le (k-1)L$
with $\rho^P_i \ge 0$, and downward the mirror image, so the slack is
$2(k-1)L$ and not $3(k-1)L$. At $k = 1$ all three terms vanish.
\end{proof}

\begin{corollary}[A work budget is sufficient]
\label{cor:sufficient}
If a policy keeps $\mathrm{In}_i - \mathrm{Out}_i \le Z$ for every job, in
particular if it keeps $\mathrm{In}_i \le Z$, then
$\mathrm{excess}_P[i] \le Z/k + (2 - 2/k)L$ for every job.
\end{corollary}

\subsection{The constant \texorpdfstring{$2(k-1)L$}{2(k-1)L} is exact}

The error term of Theorem~\ref{thm:identity} is not an artefact of the proof,
although neither search nor real input reaches it. Exhaustive
and random search over millions of small instances never exceeds $(k-1)L$ by more
than one time unit, and on the \devnum{17.6}-million-job trace of
Section~\ref{sec:experiments} the largest residual reaches \devnum{0.470} to
\devnum{0.518} of $2(k-1)L$, again about $(k-1)L$
(Section~\ref{sec:exp_bound}, Table~\ref{tab:resid}). The
extremal instances are geometric cascades needing hundreds of jobs.
% PROVENANCE (line above).  The range 0.470-0.518 is the minimum and the maximum
% of column ratio_to_bound over all 81 rows of
% outputs/dev_tables/identity_residuals.csv: min 0.46982235 (overlay 0, level 0,
% k = 8, Skip(300)), max 0.51727123 (overlay 0, level 1, k = 5, Guard(300)).
% Printed rounded outwards, so that the same 0.518 serves the "never exceeds"
% reading in Section~\ref{sec:experiments}; the stale 0.485/0.521 pair came from
% an earlier run and is not reproducible from the released package.

\begin{theorem}[Tightness of Theorem~\ref{thm:identity}]
\label{thm:tight1}
For every fixed $k \ge 2$, service bound $L>0$, and $\varepsilon > 0$ there are
inputs, work-conserving policies $P$ and jobs $i$ with
$\Delta_i > 2(k-1)L - \varepsilon$, and there are others with
$\Delta_i < -2(k-1)L + \varepsilon$. The construction is first given at scale
$\widehat L=k^m$ and then every arrival time and service time is multiplied by
$L/\widehat L$; thus the limiting sequence keeps the stated $L$ fixed. Hence
$\sup |\Delta_i| = 2(k-1)L$, a supremum that is never attained.
\end{theorem}

Both witnesses start from the cascade of Appendix~\ref{app:cascade}, which drives
$U_P - U_{\mathrm{FCFS}}$ to $(k-1)L\,(1 - ((k-1)/k)^m)$ in $m$ rounds. The
upward witness then injects $k-1$ long jobs, the victim, and a stream of unit
jobs the base policy prefers, so that first-come first-served starts the victim
at once with $\rho^{\mathrm{FCFS}}_i = (k-1)L$ while $P$ holds it until the
stream runs out; the downward witness is the mirror image, with $P$ dispatching
$k-1$ jobs of service $L$ in the victim's own phase. Appendix~\ref{app:tight1}
gives both in full, with the exact values they attain. Non-attainment follows
from Lemma~\ref{lem:gap}.

The two directions do not have the same content. Upward the extremal value is a
statement about delay, and the victim's own excess is about $+3L$; downward half of
it is the bookkeeping convention that charges an overtaker dispatched in the victim's
own phase, which executes nothing while the victim waits, and there the victim's own
excess is negative. Supplementary Section~S4.3 gives the remark in full, with the
same-phase shares, the executed-work variant of the identity, under which the upward
supremum is unchanged and the best downward witness falls to $(k-1)L$, and the two
reasons we keep the dispatch-sequence convention.



\subsection{Necessity, and What a Count of Overtakes Costs}

Theorem~\ref{thm:identity} has a converse. A policy that promises anything
about individual jobs is forced to bound overtaken work.

\begin{lemma}
\label{lem:outbound}
Let $i$ be a job, and suppose $\mathrm{excess}_P[i] \le G$ and
$\mathrm{excess}_P[j] \le G$ for every job $j$ counted in $\mathrm{Out}_i$. Then
$\mathrm{Out}_i \le k\,(G - \mathrm{excess}_P[i] + L)$.
\end{lemma}



\begin{corollary}[A work budget is necessary]
\label{cor:necessary}
If $P$ satisfies $\mathrm{excess}_P[j] \le G$ for every job on some input, then
$\mathrm{In}_i \le kG + (3k-2)L$ for every job $i$.
\end{corollary}



Lemma~\ref{lem:outbound} holds because every job counted in $\mathrm{Out}_i$ starts
inside a window of length $G - \mathrm{excess}_P[i]$, and of the jobs one server
starts inside such a window all but the last complete inside it;
Corollary~\ref{cor:necessary} then adds it to Theorem~\ref{thm:identity}, where the
terms in $\mathrm{excess}_P[i]$ cancel. Supplementary Section~S4.7 gives both proofs.

Together the two corollaries say that bounding $\mathrm{In}_i$ and offering a
per-job guarantee are the same requirement, up to an additive $O(kL)$ in the
budget. Neither says that Algorithm~\ref{alg:guard} is the least restrictive
rule with its guarantee, and we do not claim it.

\begin{remark}[Counting overtakes is not necessary, and is sufficient only at
the worst-case charge]
\label{rem:counts}
Bounded-overtaking mechanisms in the literature limit how many times a job may
be passed~\cite{grosof2021nudge,vanhoudt2022nudgek,grosof2022wcfs}. Inside the
model of Section~\ref{sec:model}, where $0 < C_i \le L$, a count bound does give
a per-job guarantee; it is not necessary for one, and it buys the guarantee at a
factor $L/\bar{C}$. Supplementary Section~S4.4 proves both halves, with the instance
outside the model on which a single overtake makes the excess arbitrarily large and
the instance inside it on which the count is unbounded while every job stays inside
its promise.
How the count is charged decides which guarantee. Charge it when an overtaker is
\emph{dispatched}, the reading a position counter in the prior art implements.
Every overtaker of $i$ is dispatched at an instant at which no job fires, since a
firing dispatch serves the queue head and the head is not an overtaker of $i$; so
at most $\kappa$ overtakers are dispatched while $i$ waits,
$\mathrm{In}_i \le \kappa L$, and
Corollary~\ref{cor:sufficient} gives
\begin{equation}
  W_P[i] \;\le\; W_{\mathrm{FCFS}}[i] + (\kappa + 2k - 2)\,L/k ,
  \label{eq:skip}
\end{equation}
which is the accounting and the formula used wherever this baseline appears
below. Inverting it, a promise $G$ admits
$\kappa = \lfloor Gk/L \rfloor - (2k-2)$ passes, the largest integer with
$(\kappa + 2k-2)L/k \le G$. Charging
the count at completion instead leaves up to $k-1$ overtakers in service and the
one being dispatched uncounted, which costs a further $L$ on the right and gives
the baseline fewer passes at the same promise. Either accounting charges every
overtaker the worst case $L$, so the rule is loose against a work budget by the
ratio $L/\bar{C}$, \devnum{one to two orders of magnitude} on our traces.
\end{remark}

\subsection{The guarantee of the guard}

We now bound what Algorithm~\ref{alg:guard} does. Nothing below restricts the
base policy, which may be randomised, learned, or clairvoyant.

\begin{theorem}[Guard bound]
\label{thm:guard}
Fix any base policy $A$ and any cap $B_{\max} \ge 0$, and run
Algorithm~\ref{alg:guard} with \emph{any} budget rule that respects the cap,
that is, with any function $\mathrm{budget}(q,t)$ satisfying
$0 \le \mathrm{budget}(q,t) \le B_{\max}$ for every waiting job $q$ and every
dispatch epoch $t$; the rule may depend on the job, may decrease over time, and
may take different values at two dispatches made at the same instant. For every
input and every job $i$, writing $W = W_P[i]$ for the wait under the wrapper,
\begin{align}
  \mathrm{In}_i &\;<\; B_{\max} + kL ,
  \label{eq:guardin}\\
  W &\;<\; W_{\mathrm{FCFS}}[i] + B_{\max}/k + (3 - 2/k)L ,
  \label{eq:guardadd}
\end{align}
where \eqref{eq:guardin} is asserted whenever $i$ has at least one overtaker.
\end{theorem}



The proof reads the last overtaker of $i$: at the instant it is dispatched, $i$ is
waiting and is not in the fired set, so $\mathrm{over}[i]$ is below its budget there;
the overtakers dispatched earlier and not yet finished occupy at most $k-1$ distinct
servers, and the last one adds a further $L$, which is the $kL$ of
\eqref{eq:guardin}. Theorem~\ref{thm:identity} with $\mathrm{Out}_i \ge 0$ turns that
into \eqref{eq:guardadd}. Supplementary Section~S4.5 gives the proof in full.

Nothing in the proof reads the budget rule except through the cap, so the shape
of the budget below the cap is a free choice and the promise is unaffected by
it. One shape carries a second bound, in which the allowance grows with the wait
the job has already had rather than being flat in it.

\begin{theorem}[Multiplicative bound for the age-relative shape]
\label{thm:guardmult}
Suppose in addition that
$\mathrm{budget}(q,t) = \min\bigl(B_0 + \eta k (t - a_q),\ B_{\max}\bigr)$ with
one constant $B_0 \ge 0$ for every job and $\eta \in [0,1)$. Then for every
input and every job $i$,
\begin{align}
  \mathrm{In}_i &\;<\; B_0 + \eta k W + kL ,
  \label{eq:guardinmult}\\
  (1 - \eta)\,W &\;<\; W_{\mathrm{FCFS}}[i] + B_0/k + (3 - 2/k)L ,
  \label{eq:guardmult}
\end{align}
with \eqref{eq:guardinmult} asserted whenever $i$ has at least one overtaker.
\end{theorem}





The same computation runs at $t \le s^P_i = a_i + W$, which bounds
$\mathrm{budget}(i,t)$ by $B_0 + \eta k W$; Supplementary Section~S4.5 gives it in
full. It goes through for a quantity fixed at $i$'s own arrival, so with the
queue-length shape of Section~\ref{sec:scheduling}, where
$B_{0,q} = B_0 + \gamma n_q$ varies from job to job, the multiplicative bound stated
with a common $B_0$ is false by a margin that grows with $\gamma$: Supplementary
Section~S4.6 gives the counterexample. For $\gamma > 0$ we therefore claim
Theorem~\ref{thm:guard} and not Theorem~\ref{thm:guardmult}. The cap, and with it
the promise, is unaffected either way.

A scheduler that learns $C_j$ at completion cannot charge work in progress, so
only completed work is charged; the $kL$ in \eqref{eq:guardin} is the price of
that and is what forces the additive term. Inverting \eqref{eq:guardadd} gives
the parameter rule of Section~\ref{sec:scheduling}:

\begin{corollary}
\label{cor:params}
$B_{\max} = k\bigl(G - (3 - 2/k)L\bigr)$ makes Algorithm~\ref{alg:guard} offer
the service guarantee $G$, whatever the budget shape below the cap and whatever
the base policy $A$. The choice is available when $G \ge (3 - 2/k)L$, and gives a positive budget only when
the inequality is strict; at equality $B_{\max} = 0$ and the wrapper is first-come
first-served.
\end{corollary}

\subsection{Both constants of the guarantee are exact}

\begin{theorem}[Tightness of Theorem~\ref{thm:guard}]
\label{thm:tight2}
Take the constant budget $\mathrm{budget}(q,t) \equiv B = B_{\max}$, and write
$c(k) = \sup (k\,\mathrm{excess}_P[i] - B)/L$ over all inputs, base policies and
budgets; by Theorem~\ref{thm:guard} the supremum over all budget rules that
respect the cap $B$ is the same number. Then
\begin{enumerate}
\item the term $B/k$ is attained: on an explicit family the excess is exactly
$\lceil B/k \rceil$, hence $B/k$ whenever $k$ divides $B$;
\item every finite instance satisfies
$(k\,\mathrm{excess}_P[i]-B)/L<3k-2$, while its supremum is exactly $3k-2$
for every $k \ge 1$. Equivalently, the additive constant $(3 - 2/k)L$ of
\eqref{eq:guardadd} cannot be lowered for any $k$.
\end{enumerate}
\end{theorem}

Part (i) is a queue of $k$ jobs of service $L$, the victim, and a stream of unit
jobs: the guard fires at $L + \lceil B/k \rceil$. For part (ii) the upper bound
is Theorem~\ref{thm:guard}, strictly; the lower bound runs the cascade for $m$
rounds and then injects a sequence timed so that $\mathrm{over}[i]$ stays below
the budget $B = f + L + 1$ at three successive dispatch epochs --- at $0$, at
$f$, and at $f + L = B - 1$, the last of them by a single unit --- letting the
base policy take every free server three times before the guard fires.
Appendix~\ref{app:tight2} gives both
constructions and the values they reach.

Both sources of slack in Theorem~\ref{thm:guard}, the constant of
Theorem~\ref{thm:identity} and the $kL$ of \eqref{eq:guardin}, are saturated
on the same instance, and the excess is real waiting time rather than bookkeeping:
every overtaker completes before the victim starts, so the same-phase share of
$\mathrm{In}_i$ is zero and the executed-work convention of
Supplementary Section~S4.3 gives the same number. The additive constant
$(3 - 2/k)L$ grows with $k$, from $L$ at $k = 1$ towards $3L$, and search does not
see this: a hill-climb over instances of at most \devnum{20} jobs suggests a flat
value near \devnum{$1.75L$}, while the witness above needs \devnum{396} jobs at
$k = 2$.

\subsection{When the Wrapper Changes Nothing}

\begin{proposition}[Consistency, one direction]
\label{prop:consistency}
Run the base policy $A$ alone. If $\mathrm{over}[q](t) < \mathrm{budget}(q,t)$
for every waiting $q$ at every dispatch epoch of that run, then the wrapped run
is identical to it, with the same dispatch sequence and the same waits.
Consequently, for a constant budget $B$, each of the following implies
$\text{guard}(A,B) = A$: $\max_i \mathrm{In}^A_i < B$; or $A$ itself satisfies
$\mathrm{excess}_A[i] \le G$ on this input and $B > kG + (3k-2)L$.
\end{proposition}

The proof inducts on dispatch epochs and is in Appendix~\ref{app:consistency},
with the two counterexamples that show the converse fails and that there is no
multiplicative consistency: the guard can fire and happen to choose what the base
would have chosen, and the ratio of mean waits between the wrapper at $B = 0$ and
a base ordering by service time is unbounded in $L$.

\subsection{Which assumptions can be dropped}

\begin{proposition}[Relaxations]
\label{prop:relax}
In the setting of Theorem~\ref{thm:identity} and Theorem~\ref{thm:guard}:
\begin{enumerate}
\item \emph{No global bound on service times is needed}, provided the prefix
maximum $\Lambda(t) = \max\{C_j : a_j \le t\}$ is placed correctly, and the
right place is not the same for every statement. The upward
direction of Theorem~\ref{thm:identity}, and hence
Corollary~\ref{cor:sufficient}, hold with $L$ replaced
by $\Lambda(a_i)$, since every job counted in $R^P_i$, $R^{\mathrm{FCFS}}_i$ or
$\rho^{\mathrm{FCFS}}_i$ has rank below $i$ and so arrived by $a_i$. The
wrapper bounds \eqref{eq:guardin} and \eqref{eq:guardadd}, and likewise
\eqref{eq:guardinmult} and \eqref{eq:guardmult}, need $\Lambda(s^P_i)$: the
$kL$ term of \eqref{eq:guardin} charges the last overtaker and the at most
$k-1$ overtakers still in service beside it, all of them jobs of rank above $i$,
which arrive after $a_i$ and are only bounded by the prefix maximum at the
instant $i$ starts. With $\Lambda(a_i)$ both are false, and by an unbounded
margin: at $k = 1$ with budget $B = 2$, let $i$ and one job of service $1$
arrive at $t = 0$ and a job of service $M$ at $t = 1$, with a base policy that
takes the two others first. Then $\Lambda(a_i) = 1$ for every $M$, the guard
never fires because $\mathrm{over}[i]$ stops at $1 < 2$, and
$\mathrm{In}_i = M + 1$ against a claimed $B + k\Lambda(a_i) = 3$. Written out,
the sharp forms are
$\mathrm{In}_i < B_{\max} + k\,\Lambda(s^P_i)$ and
$\mathrm{excess}_P[i] < B_{\max}/k + \Lambda(s^P_i)
 + (2-2/k)\,\Lambda(a_i)$,
and since $\Lambda(a_i) \le \Lambda(s^P_i)$ they give
\eqref{eq:guardadd} with $\Lambda(s^P_i)$ throughout. The downward
direction of Theorem~\ref{thm:identity} holds with
$\Lambda(\max(s^P_i, s^{\mathrm{FCFS}}_i))$, which we do not claim to be the
smallest horizon that works, and Corollary~\ref{cor:necessary}
holds with $\Lambda(s^{\mathrm{FCFS}}_i + G)$, since Lemma~\ref{lem:outbound}
bounds jobs that start as late as that; both are false with $\Lambda(a_i)$.
\item \emph{Servers may have different speeds.} With speeds $v_1, \dots, v_k$
and $V = \sum_{\ell} v_{\ell}$, Theorem~\ref{thm:identity} holds with $k$ replaced by $V$
on the left and the same $2(k-1)L$ on the right, and the guard satisfies
$W_P[i] \le W_{\mathrm{FCFS}}[i] + B/V + (3k-2)L/V$: the budget is divided by
total capacity, not by the server count.
Here $L$ bounds a job's work in a common unit, not its wall-clock duration: a common
wall-clock limit does not produce policy-independent executed work when speeds
differ. The reference assigns each dispatched job to the fastest free server, ties
by server index.
\item \emph{The timeout is already in the model.} Section~\ref{sec:model} defines
$C_j$ as executed work, $\min(C_j, L)$ for a job stopped at the limit, so every
currency above ($\mathrm{In}_i$, $\mathrm{Out}_i$, $\mathrm{over}[q](t)$ and
$L$) is in executed work and no assumption on raw service demands is needed.
Read in raw work the statements are false, a stopped job contributing more to
$\mathrm{In}_i$ than is ever executed.
\item \emph{Setup times survive.} If a setup of length $\theta$ is paid before
every service under both policies, absorb it into the job as
$C'_j = C_j + \theta$; the model is unchanged and every result holds with
$L \mapsto L + \theta$.
\item \emph{Work conservation cannot be dropped.} If servers may idle while jobs
wait, let $\Gamma^P_i$ be the server-idle time accumulated during
$[a_i, s^P_i)$ while at least one job waited. Then \eqref{eq:busy} becomes
$k W_P[i] = R^P_i + \mathrm{In}_i - \mathrm{Out}_i - \rho^P_i + \Gamma^P_i$, and
nothing survives without a bound on $\Gamma$: with $k = 1$, pausing the server
for $T$ while $i$ waits leaves $\mathrm{In}_i = \mathrm{Out}_i = 0$ and makes
the excess $T$. Release constraints are the same phenomenon: if rank is keyed
on arrival while eligibility starts later, the identity fails by exactly the
forced idle time, and ranking on release repairs it.
\end{enumerate}
\end{proposition}

\subsection{A floor of order \texorpdfstring{$L$}{L}, and what a guarantee costs}

\begin{proposition}[Impossibility]
\label{prop:floor}
Fix any budget rule that assigns a positive budget to every waiting job at the
first dispatch epoch. Then there are a base policy $A$ that observes a job's
service time only at its completion, and an input, such that
Algorithm~\ref{alg:guard} run around $A$ with that rule leaves some job with
$\mathrm{excess} \ge L$. For a base policy that ranks by predictions, such as
SPJF-E, the predictions in the input are enough and $A$ need not be chosen.
Consequently no promise $G < L$ can be made by a wrapper of this family that
holds for every base policy, or
for one prediction-ranked base policy under arbitrary predictions, which is the
generality in which Theorem~\ref{thm:guard} is stated. At $k=1$ and constant
budget $B>0$, the bound $B+L$ of Theorem~\ref{thm:guard} is the exact supremum
of the excess, approached and not attained; at $B=0$ the wrapper is FCFS and
the excess is exactly zero.
\end{proposition}




The proof puts $k+1$ jobs of service $L$ into an empty system at $t = 0$: nothing has
completed, so no waiting job fires, the wrapper accepts the base policy's choices, and
the first choice beyond the first $k$ ranks displaces a job that first-come
first-served starts at once, leaving it with excess $L$. The hypothesis can be neither
dropped nor tightened: as a claim about policies the statement is false, since
first-come first-served itself has excess $0$ on every input, and a rule that is $0$ at
the first dispatch epoch and positive later falls to the same instance placed behind a
saturating prefix. Supplementary Section~S4.8 gives the proof and the case analysis
that makes those cases exhaustive for the budget of Algorithm~\ref{alg:guard}, which
either makes the wrapper first-come first-served around every base policy or admits a
base policy and an input with excess at least $L$.

The guarantee also caps how much reordering can buy at all, and not only how much
it may cost a single job.

\begin{theorem}[Price of the guarantee]
\label{thm:price}
Take $m$ jobs of service $L$ and $n$ jobs of service $\sigma < L$, the long ones
ranked first, all arriving at $t = 0$, and let $P$ be any non-preemptive
work-conserving policy with $\mathrm{excess}_P[j] \le G$ for every job. Define
$D_0=mn(L-\sigma)/k$, $E=2(1-1/k)L(m+n)$ and
$A=(kG+(3k-2)L)/(n\sigma)$. If
$D_0>E$, equivalently
$mn(L-\sigma)>2(k-1)L(m+n)$, then
\[
  \frac{\overline{W}_{\mathrm{FCFS}} - \overline{W}_{P}}
       {\overline{W}_{\mathrm{FCFS}} - \overline{W}_{\mathrm{SJF}}}
  \;\le\; A+\varepsilon_+,
  \qquad
  \varepsilon_+=\frac{E(1+A)}{D_0-E}.
\]
For fixed $k,L,\sigma,G$, the finite-sample term is
$O((m+n)/(mn))$.
\end{theorem}

The proof is in Appendix~\ref{app:price}. The right-hand side is capacity times
guarantee, divided by the work that wants to overtake; it is linear in $G$ and
saturates at $1$ when $kG$ reaches the total service of the impatient class.
Because it applies to every policy with the guarantee and not only to
Algorithm~\ref{alg:guard}, it describes the trade-off an operator faces rather
than a weakness of this wrapper. The additive $(3k-2)L$ in it is slack: it comes
from applying Corollary~\ref{cor:necessary} to each long job separately, where
the long jobs can be charged together, and at $k = 1$ the aggregate argument of
Appendix~\ref{app:price} removes it. There the achievable fraction is known
exactly: it is $\lfloor G/\sigma \rfloor / n$, capped at $1$
(Proposition~\ref{prop:k1frac}).

\subsection{Achievability and maximal permission}

The preceding results could otherwise be read only as a conjunction of a robust
wrapper and a necessary budget. The next theorem matches the price bound, up to
its stated finite-sample band, on the two-size batch family; the theorem after it
identifies the pointwise-maximal permission rule in a precisely restricted
single-server class. Neither result gives instance optimality on arbitrary
inputs, and the exact value of the optimum below remains open for $k\geq2$.

\begin{theorem}[Achievability on the two-size batch family]
\label{thm:two_size_achievability}
In the family of Theorem~\ref{thm:price}, define
\[
  \Delta(Q)=\sum_i\bigl(W_{\mathrm{FCFS}}[i]-W_Q[i]\bigr),\qquad
  F(Q)=\frac{\Delta(Q)}{\Delta(\mathrm{SJF})},
\]
and let $F^*(G)$ be the maximum of $F(P)$ over all non-preemptive,
work-conserving policies satisfying $\mathrm{excess}_P[i]\leq G$ for every job.
Let the guard use exact SJF as its base and the constant budget
$B=kG-(3k-2)L\geq0$. Put
\[
 D_0=\frac{mn(L-\sigma)}{k},\quad
 E=2(1-1/k)L(m+n),\quad
 q=\min\!\left(1,\frac{B}{n\sigma}\right),\quad
 A=\frac{kG+(3k-2)L}{n\sigma}.
\]
If $D_0>E$, then
\[
 q-\varepsilon_-\ \leq\ F(\mathrm{guard})\ \leq\ F^*(G)
 \ \leq\ \min(1,A+\varepsilon_+),
\]
where
\[
 \varepsilon_- = \frac{E(1+q)}{D_0-E},
 \qquad
 \varepsilon_+ = \frac{E(1+A)}{D_0-E}.
\]
Thus the leading terms are separated by at most
$2(3k-2)L/(n\sigma)$, in addition to the two displayed finite-sample
errors. For $k\geq2$, this sandwich does not determine the exact $F^*(G)$.
\end{theorem}

\begin{definition}[Opaque-base, size-oblivious wrapper]
\label{def:opaque_wrapper}
At a dispatch epoch, the wrapper may use arrivals and ranks, its own dispatch and
completion history, the revealed sizes of completed jobs, the public bound $L$,
fixed random coins, and the sequence of proposals made so far by an opaque base
policy. It may not inspect unfinished job sizes or the base policy's internal
state. A deterministic pathwise promise $G$ means that every job has excess at
most $G$ for every finite input and every possible base-proposal transcript.
\end{definition}

\begin{theorem}[Pointwise maximality at one server]
\label{thm:k1_maximal}
Assume $k=1$, sizes and budgets are nonnegative integer time units bounded by the
positive integer $L$, and $G\geq L$ is an integer. Within
Definition~\ref{def:opaque_wrapper}:
\begin{enumerate}
\item for constant budget $B\geq1$, the guard's worst pathwise excess is exactly
$B+L-1$; at $B=0$ it is FCFS and the excess is zero;
\item the largest positive integer constant budget with promise $G$ is
$B^*=G-L+1$;
\item at any reachable history, if the queue head $h$ satisfies
$\mathrm{over}[h]>G-L$, every wrapper with promise $G$ must dispatch $h$.
If $\mathrm{over}[h]\leq G-L$, permitting any currently waiting later-ranked
job for this dispatch is safe against the worst unrevealed size $L$.
Consequently the guard triggered by $\mathrm{over}\geq B^*$ is pointwise maximal
in its permitted action set.
\end{enumerate}
\end{theorem}

\begin{corollary}[The single-server consistency--robustness frontier]
\label{cor:k1_frontier}
On the integer two-size batch family, for $G\geq L$ the pointwise-maximal guard
around the exact-SJF proposal sequence closes the fraction
\[
 \min\!\left(1,\frac{\lfloor(G-L)/\sigma\rfloor+1}{n}\right)
\]
of the FCFS-to-SJF total-wait gap, and no opaque-base wrapper with the pathwise
promise can permit a larger fraction of those short-before-long moves. Exact SJF
itself reads unfinished true sizes, so the composed exact-SJF-plus-guard scheduler
is not size-oblivious; the frontier concerns the wrapper's permissions under an
externally supplied proposal transcript.
\end{corollary}

Appendix~\ref{app:achievability} proves these statements. Supplementary theory
material records an unbounded instance-optimality counterexample and the limits
of instance-dependent refinements.

Every uniform constant above is on the scale of $L$. Proposition~\ref{prop:floor}
establishes such a floor for the wrapper family under arbitrary opaque base
proposals, not for all scheduling mechanisms. Empirically, usefulness depends on
size dispersion, score predictability and the arrival structure; the CI pools
are negative cases rather than evidence of a necessary distributional condition.
Section~\ref{sec:exp_cross} quantifies that boundary on the evaluated workloads.
